\documentclass[11pt,a4paper]{article}
\usepackage[T1]{fontenc}
\usepackage[utf8]{inputenc}
\usepackage[brazil]{babel}
\usepackage{amsmath,amssymb,textcomp}
\usepackage[margin=2.2cm]{geometry}
\usepackage{enumitem}
\usepackage{array}
\setlength{\parindent}{0pt}
\setlength{\parskip}{0.6em}

\begin{document}
\section*{Avaliação de questões --- avaliador 5}
\subsection*{Como responder}

São 21 questões de Matemática do Ensino Médio. Leia cada uma como leria uma
questão que você pensasse em usar com a sua turma, e responda às quatro
perguntas que vêm logo abaixo dela.

Onde se pergunta pelo \textbf{nível de dificuldade}, responda pensando nos
seus alunos, não num aluno ideal. É essa resposta que permite calibrar o
sistema: ele classifica cada questão como fácil, média ou difícil por conta
própria, e precisamos saber o quanto isso corresponde à sala de aula real.

\textbf{A amostra tem qualidade variável, e pode conter questões com erro
matemático.} Se encontrar algum, aponte onde.

Não há resposta certa sobre você: o que está sendo avaliado é o sistema que
produziu estas questões, não quem as lê. Um ``recusada'' bem justificado vale
mais para este trabalho do que um ``aceita'' por gentileza.

Você pode responder direto neste documento impresso ou na planilha
\texttt{respostas-avaliador-5.csv}, que já vem com os códigos dos itens nas linhas.

Tempo estimado: cerca de 75 minutos.
\bigskip\hrule\bigskip
\begin{samepage}
\subsection*{Q007}
Uma oficina mecânica cobra por seus serviços de revisão um valor fixo de mão de obra, mais uma taxa que depende do número de horas de trabalho. Uma revisão que durou 2 horas custou R\$\,150,00, e outra revisão, no mesmo padrão de cobrança, durou 5 horas e custou R\$\,300,00. Sabendo que o custo total varia linearmente com o número de horas trabalhadas, qual seria o custo de uma revisão que durasse 8 horas?
\begin{enumerate}[label=(\alph*),nosep,leftmargin=*]
\item R\$\,450,00
\item R\$\,400,00
\item R\$\,600,00
\item R\$\,480,00
\end{enumerate}

\textbf{Gabarito proposto:} R\$\,450,00

\textbf{Habilidade declarada:} EM13MAT302 --- Resolver e elaborar problemas cujos modelos são as funções polinomiais de 1º e 2º graus, em contextos diversos, incluindo ou não tecnologias digitais.
\end{samepage}

\noindent\rule{\linewidth}{0.4pt}
\textbf{Tem erro matemático?}\quad
\mbox{$\square$~não} \quad \mbox{$\square$~sim, aqui:} \underline{\hspace{5cm}}
\quad \mbox{$\square$~não sei dizer}

\textbf{Que nível de dificuldade esta questão tem para os seus alunos?}\quad
\mbox{$\square$~fácil} \quad \mbox{$\square$~média} \quad
\mbox{$\square$~difícil} \quad \mbox{$\square$~fora do alcance da turma}

\textbf{Você usaria esta questão?}\quad
\mbox{$\square$~aceita} \quad \mbox{$\square$~aceita com ajuste} \quad
\mbox{$\square$~recusada}

\textbf{Por quê?} \emph{(obrigatório quando não for ``aceita'')}

\underline{\hspace{\linewidth}}

\underline{\hspace{\linewidth}}
\vspace{0.6em}

\begin{samepage}
\subsection*{Q024}
A tabela abaixo relaciona valores de $x$ e os correspondentes valores de $y$, obtidos a partir de uma mesma regra de formação (os valores de $x$ não variam em intervalos iguais):

\begin{center}
\begin{tabular}{|p{0.172\linewidth}|p{0.172\linewidth}|p{0.172\linewidth}|p{0.172\linewidth}|p{0.172\linewidth}|}
\hline
\textbf{$x$} & \textbf{$-3$} & \textbf{$1$} & \textbf{$4$} & \textbf{$8$} \\
\hline
$y$ & $14$ & $2$ & $-7$ & $-19$ \\
\hline
\end{tabular}
\end{center}

Um estudante, analisando os pares ordenados $(x,y)$ da tabela, conjectura que essa relação pode ser representada por uma função polinomial do 1º grau. Assinale a alternativa que apresenta a lei de formação $f(x)$ que generaliza corretamente o padrão observado na tabela.
\begin{enumerate}[label=(\alph*),nosep,leftmargin=*]
\item $f(x) = -3x + 5$
\item $f(x) = -12x + 14$
\item $f(x) = 3x - 1$
\item $f(x) = -11x + 13$
\end{enumerate}

\textbf{Gabarito proposto:} f(x) = -3x + 5

\textbf{Habilidade declarada:} EM13MAT501 --- Investigar relações entre números expressos em tabelas para representá-los no plano cartesiano, identificando padrões e criando conjecturas para generalizar e expressar algebricamente essa generalização, reconhecendo quando essa representação é de função polinomial de 1º grau.
\end{samepage}

\noindent\rule{\linewidth}{0.4pt}
\textbf{Tem erro matemático?}\quad
\mbox{$\square$~não} \quad \mbox{$\square$~sim, aqui:} \underline{\hspace{5cm}}
\quad \mbox{$\square$~não sei dizer}

\textbf{Que nível de dificuldade esta questão tem para os seus alunos?}\quad
\mbox{$\square$~fácil} \quad \mbox{$\square$~média} \quad
\mbox{$\square$~difícil} \quad \mbox{$\square$~fora do alcance da turma}

\textbf{Você usaria esta questão?}\quad
\mbox{$\square$~aceita} \quad \mbox{$\square$~aceita com ajuste} \quad
\mbox{$\square$~recusada}

\textbf{Por quê?} \emph{(obrigatório quando não for ``aceita'')}

\underline{\hspace{\linewidth}}

\underline{\hspace{\linewidth}}
\vspace{0.6em}

\begin{samepage}
\subsection*{Q073}
O nível da maré em certo ponto da costa varia de forma cíclica ao longo do dia e pode ser modelado pela função $h(t) = 2\cos\left(\dfrac{\pi}{6}(t-3)\right) + 3$, em que $t$ é o tempo em horas contado a partir da meia-noite ($t \ge 0$) e $h(t)$ é a altura da maré, em metros. Interpretando o gráfico dessa função como uma curva cossenoidal no plano cartesiano, determine: a altura máxima e a altura mínima da maré, o período do fenômeno e todos os instantes $t$, com $0 \le t < 12$, em que a maré atinge exatamente 4 metros de altura.
\begin{enumerate}[label=(\alph*),nosep,leftmargin=*]
\item Altura máxima 5 m; altura mínima 1 m; período 12 h; maré de 4 m em $t=1$ h e $t=5$ h.
\item Altura máxima 2 m; altura mínima -2 m; período 12 h; maré de 4 m em $t=1$ h e $t=5$ h.
\item Altura máxima 5 m; altura mínima 1 m; período 6 h; maré de 4 m em $t=1$ h e $t=5$ h.
\item Altura máxima 5 m; altura mínima 1 m; período 12 h; maré de 4 m apenas em $t=1$ h.
\end{enumerate}

\textbf{Gabarito proposto:} Altura máxima 5 m; altura mínima 1 m; período 12 h; maré de 4 m em $t=1$ h e $t=5$ h.

\textbf{Habilidade declarada:} EM13MAT306 --- Resolver e elaborar problemas em contextos que envolvem fenômenos periódicos reais, como ondas sonoras, ciclos menstruais, movimentos cíclicos, entre outros, e comparar suas representações com as funções seno e cosseno, no plano cartesiano, com ou sem apoio de aplicativos de álgebra e geometria.
\end{samepage}

\noindent\rule{\linewidth}{0.4pt}
\textbf{Tem erro matemático?}\quad
\mbox{$\square$~não} \quad \mbox{$\square$~sim, aqui:} \underline{\hspace{5cm}}
\quad \mbox{$\square$~não sei dizer}

\textbf{Que nível de dificuldade esta questão tem para os seus alunos?}\quad
\mbox{$\square$~fácil} \quad \mbox{$\square$~média} \quad
\mbox{$\square$~difícil} \quad \mbox{$\square$~fora do alcance da turma}

\textbf{Você usaria esta questão?}\quad
\mbox{$\square$~aceita} \quad \mbox{$\square$~aceita com ajuste} \quad
\mbox{$\square$~recusada}

\textbf{Por quê?} \emph{(obrigatório quando não for ``aceita'')}

\underline{\hspace{\linewidth}}

\underline{\hspace{\linewidth}}
\vspace{0.6em}

\begin{samepage}
\subsection*{Q065}
Um oceanógrafo registra a altura da maré (em metros) em um pequeno porto ao longo de um dia. Ele observa que a maré varia de forma periódica e simétrica em torno de um nível médio, atingindo seu valor máximo de 1,8 m exatamente à meia-noite ($t=0$, com $t$ em horas) e seu valor mínimo de 0,2 m exatamente às 6h da manhã. Esse padrão se repete a cada 12 horas.

a) Modele a altura da maré, em metros, por uma função do tipo $h(t) = A\cos(Bt) + D$, com $t$ em horas contado a partir da meia-noite. Determine explicitamente os valores de $A$, $B$ e $D$, justificando cada um a partir dos dados do fenômeno (amplitude, período e deslocamento vertical do gráfico em relação ao eixo $t$).

b) Determine todos os horários dentro das primeiras 12 horas do dia (isto é, para $0 \le t < 12$) em que a maré atinge exatamente 1,4 m de altura. Para cada horário encontrado, indique se a maré está subindo ou descendo naquele instante, justificando sua resposta apenas a partir do comportamento gráfico da função cosseno (sem usar derivadas).

\textbf{Gabarito proposto:} h(t) = 0,8·cos($\pi$t/6) + 1,0 (A=0,8 m; B=$\pi$/6 rad/h; D=1,0 m). A maré atinge 1,4 m às 2h (descendo) e às 10h (subindo).

\textbf{Habilidade declarada:} EM13MAT306 --- Resolver e elaborar problemas em contextos que envolvem fenômenos periódicos reais, como ondas sonoras, ciclos menstruais, movimentos cíclicos, entre outros, e comparar suas representações com as funções seno e cosseno, no plano cartesiano, com ou sem apoio de aplicativos de álgebra e geometria.
\end{samepage}

\noindent\rule{\linewidth}{0.4pt}
\textbf{Tem erro matemático?}\quad
\mbox{$\square$~não} \quad \mbox{$\square$~sim, aqui:} \underline{\hspace{5cm}}
\quad \mbox{$\square$~não sei dizer}

\textbf{Que nível de dificuldade esta questão tem para os seus alunos?}\quad
\mbox{$\square$~fácil} \quad \mbox{$\square$~média} \quad
\mbox{$\square$~difícil} \quad \mbox{$\square$~fora do alcance da turma}

\textbf{Você usaria esta questão?}\quad
\mbox{$\square$~aceita} \quad \mbox{$\square$~aceita com ajuste} \quad
\mbox{$\square$~recusada}

\textbf{Por quê?} \emph{(obrigatório quando não for ``aceita'')}

\underline{\hspace{\linewidth}}

\underline{\hspace{\linewidth}}
\vspace{0.6em}

\begin{samepage}
\subsection*{Q084}
Dois estudantes lançam, ao mesmo tempo, foguetes de brinquedo verticalmente para cima a partir do solo. As alturas atingidas pelos foguetes A e B, em metros, em função do tempo $t$ (em segundos) após o lançamento, são dadas por:

$h_A(t) = -5t^2 + 20t + 2$

$h_B(t) = -4t^2 + 18t + 3$

Qual das alternativas a seguir indica corretamente qual foguete atinge a maior altura máxima, o valor dessa altura, o instante em que ela ocorre, e justifica por que esse valor corresponde de fato a um ponto de máximo (e não de mínimo)?
\begin{enumerate}[label=(\alph*),nosep,leftmargin=*]
\item O foguete B atinge a maior altura máxima, 23,25 m, no instante t = 2,25 s; como o coeficiente de $t^2$ é negativo em ambas as funções, seus gráficos têm concavidade voltada para baixo e o vértice é ponto de máximo.
\item O foguete A atinge a maior altura máxima, pois seu coeficiente de $t^2$ (-5) é mais negativo que o do foguete B (-4), o que indicaria uma parábola 'mais alta'.
\item O foguete B atinge a maior altura máxima, 20,25 m, no instante t = 2,25 s.
\item O foguete B atinge a maior altura máxima, pois sua altura inicial (3 m) é maior que a do foguete A (2 m).
\end{enumerate}

\textbf{Gabarito proposto:} O foguete B atinge a maior altura máxima, 23,25 m, no instante t = 2,25 s; como o coeficiente de t² é negativo em ambas as funções, seus gráficos têm concavidade voltada para baixo e o vértice é ponto de máximo.

\textbf{Habilidade declarada:} EM13MAT503 --- Investigar pontos de máximo ou de mínimo de funções quadráticas em contextos da Matemática Financeira ou da Cinemática, entre outros.
\end{samepage}

\noindent\rule{\linewidth}{0.4pt}
\textbf{Tem erro matemático?}\quad
\mbox{$\square$~não} \quad \mbox{$\square$~sim, aqui:} \underline{\hspace{5cm}}
\quad \mbox{$\square$~não sei dizer}

\textbf{Que nível de dificuldade esta questão tem para os seus alunos?}\quad
\mbox{$\square$~fácil} \quad \mbox{$\square$~média} \quad
\mbox{$\square$~difícil} \quad \mbox{$\square$~fora do alcance da turma}

\textbf{Você usaria esta questão?}\quad
\mbox{$\square$~aceita} \quad \mbox{$\square$~aceita com ajuste} \quad
\mbox{$\square$~recusada}

\textbf{Por quê?} \emph{(obrigatório quando não for ``aceita'')}

\underline{\hspace{\linewidth}}

\underline{\hspace{\linewidth}}
\vspace{0.6em}

\begin{samepage}
\subsection*{Q017}
Um biólogo estuda o crescimento de uma colônia de bactérias em laboratório. Ele constata que, em um modelo matemático simplificado, a quantidade de indivíduos na colônia, em milhares, pode ser descrita pela função $f(x) = 2^{x}$, em que $x$ representa o número de horas decorridas desde um instante de referência (podendo $x$ assumir valores negativos, representando horas anteriores a esse instante, dentro do intervalo em que o modelo é válido).

Para resolver o problema inverso — estimar quantas horas se passaram até que a colônia atingisse uma determinada quantidade $y$ de milhares de indivíduos — o biólogo utiliza a função $g(x) = \log_{2}(x)$.

Com base nessas duas funções, resolva:

\textbf{a)} Determine o domínio e a imagem de $f$ e de $g$. Compare os dois pares de conjuntos e explique, em termos matemáticos, por que essa relação entre eles não é uma coincidência.

\textbf{b)} Analise o crescimento de $f$ e de $g$ (cada uma é crescente ou decrescente?). Embora ambas cresçam no mesmo sentido, explique, sem usar tabelas numéricas, por que o crescimento de $f$ é qualitativamente muito mais rápido que o de $g$ à medida que $x$ aumenta.

\textbf{c)} Calcule as composições $f(g(x))$ e $g(f(x))$, indicando para quais valores de $x$ cada composição está definida. O que esses resultados revelam sobre a relação entre as funções $f$ e $g$?

\textbf{d)} Sabendo que $f(1) = 2$ e $f(3) = 8$, calcule $g(2)$ e $g(8)$. Em seguida, esboce em um mesmo plano cartesiano os gráficos de $f$, de $g$ e da reta $y = x$, situando nele os pontos $(1,2)$, $(2,1)$, $(3,8)$ e $(8,3)$. Explique, a partir desse esboço, que relação geométrica existe entre o gráfico de $f$ e o gráfico de $g$, e como as respostas dos itens (a), (b) e (c) justificam essa relação.

\textbf{Gabarito proposto:} a) $D_f=\mathbb{R}$, $Im_f=(0,+\infty)$; $D_g=(0,+\infty)$, $Im_g=\mathbb{R}$ — domínio e imagem se invertem entre f e g. b) f e g são ambas crescentes, mas f cresce de forma multiplicativa (cada vez mais rápido) e g cresce de forma cada vez mais lenta. c) $f(g(x))=x$ para $x>0$ e $g(f(x))=x$ para todo x real, mostrando que g é a função inversa de f. d) $g(2)=1$, $g(8)=3$; os gráficos de f e g são simétricos em relação à reta $y=x$.

\textbf{Habilidade declarada:} EM13MAT403 --- Comparar e analisar as representações, em plano cartesiano, das funções exponencial e logarítmica para identificar as características fundamentais (domínio, imagem, crescimento) de cada uma, com ou sem apoio de tecnologias digitais, estabelecendo relações entre elas.
\end{samepage}

\noindent\rule{\linewidth}{0.4pt}
\textbf{Tem erro matemático?}\quad
\mbox{$\square$~não} \quad \mbox{$\square$~sim, aqui:} \underline{\hspace{5cm}}
\quad \mbox{$\square$~não sei dizer}

\textbf{Que nível de dificuldade esta questão tem para os seus alunos?}\quad
\mbox{$\square$~fácil} \quad \mbox{$\square$~média} \quad
\mbox{$\square$~difícil} \quad \mbox{$\square$~fora do alcance da turma}

\textbf{Você usaria esta questão?}\quad
\mbox{$\square$~aceita} \quad \mbox{$\square$~aceita com ajuste} \quad
\mbox{$\square$~recusada}

\textbf{Por quê?} \emph{(obrigatório quando não for ``aceita'')}

\underline{\hspace{\linewidth}}

\underline{\hspace{\linewidth}}
\vspace{0.6em}

\begin{samepage}
\subsection*{Q087}
Considere a progressão aritmética $(a_n)$ com primeiro termo $a_1 = 7$ e razão $r = 4$: $(7, 11, 15, 19, 23, \dots)$. Essa PA pode ser vista como a restrição da função afim $f(x) = 4x + 3$ ao conjunto dos números naturais não nulos, de modo que $a_n = f(n)$ para todo $n \geq 1$ (o gráfico da PA é o conjunto de pontos discretos $(n, a_n)$ sobre a reta que representa $f$).

Qual é o menor termo dessa PA que é maior do que 100?
\begin{enumerate}[label=(\alph*),nosep,leftmargin=*]
\item 103
\item 99
\item 107
\item 100
\end{enumerate}

\textbf{Gabarito proposto:} 103

\textbf{Habilidade declarada:} EM13MAT507 --- Identificar e associar sequências numéricas (PA) a funções afins de domínios discretos para análise de propriedades, incluindo dedução de algumas fórmulas e resolução de problemas.
\end{samepage}

\noindent\rule{\linewidth}{0.4pt}
\textbf{Tem erro matemático?}\quad
\mbox{$\square$~não} \quad \mbox{$\square$~sim, aqui:} \underline{\hspace{5cm}}
\quad \mbox{$\square$~não sei dizer}

\textbf{Que nível de dificuldade esta questão tem para os seus alunos?}\quad
\mbox{$\square$~fácil} \quad \mbox{$\square$~média} \quad
\mbox{$\square$~difícil} \quad \mbox{$\square$~fora do alcance da turma}

\textbf{Você usaria esta questão?}\quad
\mbox{$\square$~aceita} \quad \mbox{$\square$~aceita com ajuste} \quad
\mbox{$\square$~recusada}

\textbf{Por quê?} \emph{(obrigatório quando não for ``aceita'')}

\underline{\hspace{\linewidth}}

\underline{\hspace{\linewidth}}
\vspace{0.6em}

\begin{samepage}
\subsection*{Q054}
Um biólogo monitora o crescimento de uma cultura de bactérias em laboratório. No início da observação (t = 0), há 400 bactérias, e a cultura cresce de modo que a população dobra a cada 6 horas. A população N(t), em número de bactérias, t horas após o início da observação, é dada por N(t) = 400 · 2\textasciicircum{}(t/6). Considerando essa lei de crescimento, qual das afirmações a seguir compara corretamente o crescimento da população entre t = 0h e t = 6h com o crescimento entre t = 12h e t = 18h?
\begin{enumerate}[label=(\alph*),nosep,leftmargin=*]
\item Em ambos os intervalos a população dobra (fator multiplicativo 2), mas o aumento absoluto de bactérias é maior no segundo intervalo: 400 entre 0h e 6h, contra 1600 entre 12h e 18h, pois a mesma taxa de crescimento multiplica quantidades cada vez maiores.
\item Em ambos os intervalos a população aumenta exatamente 400 bactérias, já que a taxa de crescimento de 2 vezes a cada 6 horas corresponde a um acréscimo fixo de indivíduos por hora.
\item No segundo intervalo (12h a 18h) a população cresce mais do que o dobro, pois quanto mais tempo passa, maior fica o fator de multiplicação aplicado a cada período de 6 horas.
\item A população dobra apenas uma vez em todo o período analisado, atingindo 1600 bactérias em 18 horas, pois o enunciado indica que o dobro ocorre a cada 12 horas, não a cada 6.
\end{enumerate}

\textbf{Gabarito proposto:} Em ambos os intervalos a população dobra (fator 2), mas o aumento absoluto é de 400 bactérias no primeiro intervalo e de 1600 bactérias no segundo, pois a mesma razão de crescimento multiplica quantidades cada vez maiores.

\textbf{Habilidade declarada:} EM13MAT304 --- Resolver e elaborar problemas com funções exponenciais nos quais é necessário compreender e interpretar a variação das grandezas envolvidas, em contextos como o da Matemática Financeira e o do crescimento de seres vivos microscópicos, entre outros.
\end{samepage}

\noindent\rule{\linewidth}{0.4pt}
\textbf{Tem erro matemático?}\quad
\mbox{$\square$~não} \quad \mbox{$\square$~sim, aqui:} \underline{\hspace{5cm}}
\quad \mbox{$\square$~não sei dizer}

\textbf{Que nível de dificuldade esta questão tem para os seus alunos?}\quad
\mbox{$\square$~fácil} \quad \mbox{$\square$~média} \quad
\mbox{$\square$~difícil} \quad \mbox{$\square$~fora do alcance da turma}

\textbf{Você usaria esta questão?}\quad
\mbox{$\square$~aceita} \quad \mbox{$\square$~aceita com ajuste} \quad
\mbox{$\square$~recusada}

\textbf{Por quê?} \emph{(obrigatório quando não for ``aceita'')}

\underline{\hspace{\linewidth}}

\underline{\hspace{\linewidth}}
\vspace{0.6em}

\begin{samepage}
\subsection*{Q026}
Uma roda-gigante tem raio de 10 m e seu eixo central está a 12 m de altura em relação ao solo. Ela gira em velocidade constante, completando uma volta a cada 40 segundos. No instante $t = 0$ (em segundos), um determinado carrinho está exatamente no ponto mais baixo do movimento circular. A altura $h(t)$, em metros, desse carrinho em relação ao solo, $t$ segundos após o início da contagem do tempo, é dada por:

$$h(t) = 12 - 10\cos\left(\dfrac{\pi t}{20}\right)$$

Imagine o movimento do carrinho representado simultaneamente de duas formas: como um ponto percorrendo o ciclo trigonométrico (com ângulo variando de $0$ a $2\pi$ radianos e voltando a se repetir) e como uma curva no plano cartesiano, com o tempo $t$ no eixo horizontal e a altura $h(t)$ no eixo vertical.

Com base nessa comparação entre as duas representações, responda:

a) Qual é o período da função $h$, ou seja, o tempo necessário para que o carrinho complete uma volta e a altura volte a se repetir? Explique como esse valor se relaciona com o fato de o ciclo trigonométrico ter $2\pi$ radianos.

b) Qual é o domínio da função $h$, considerando que $t$ representa o tempo decorrido desde o início do funcionamento da roda-gigante (que continua girando indefinidamente)?

c) Qual é a imagem da função $h$, ou seja, quais são as menores e maiores alturas atingidas pelo carrinho? Relacione essas alturas com os pontos do ciclo trigonométrico em que $\cos(\theta) = 1$ e $\cos(\theta) = -1$.

\textbf{Gabarito proposto:} a) Período $T = 40$ s, pois o ângulo $\frac{\pi t}{20}$ percorre $2\pi$ radianos (uma volta completa no ciclo trigonométrico) exatamente quando $t=40$. b) Domínio: $[0,+\infty)$. c) Imagem: $[2,22]$ metros, correspondendo a $\cos(\theta)=1$ (altura mínima, 2 m) e $\cos(\theta)=-1$ (altura máxima, 22 m).

\textbf{Habilidade declarada:} EM13MAT404 --- Identificar as características fundamentais das funções seno e cosseno (periodicidade, domínio, imagem), por meio da comparação das representações em ciclos trigonométricos e em planos cartesianos, com ou sem apoio de tecnologias digitais.
\end{samepage}

\noindent\rule{\linewidth}{0.4pt}
\textbf{Tem erro matemático?}\quad
\mbox{$\square$~não} \quad \mbox{$\square$~sim, aqui:} \underline{\hspace{5cm}}
\quad \mbox{$\square$~não sei dizer}

\textbf{Que nível de dificuldade esta questão tem para os seus alunos?}\quad
\mbox{$\square$~fácil} \quad \mbox{$\square$~média} \quad
\mbox{$\square$~difícil} \quad \mbox{$\square$~fora do alcance da turma}

\textbf{Você usaria esta questão?}\quad
\mbox{$\square$~aceita} \quad \mbox{$\square$~aceita com ajuste} \quad
\mbox{$\square$~recusada}

\textbf{Por quê?} \emph{(obrigatório quando não for ``aceita'')}

\underline{\hspace{\linewidth}}

\underline{\hspace{\linewidth}}
\vspace{0.6em}

\begin{samepage}
\subsection*{Q014}
Uma locadora de bicicletas cobra uma taxa fixa de utilização mais um valor proporcional ao tempo de uso. A tabela a seguir mostra o valor total pago por diferentes clientes, de acordo com o tempo de uso da bicicleta:

\begin{center}
\begin{tabular}{|p{0.430\linewidth}|p{0.430\linewidth}|}
\hline
\textbf{Tempo de uso (horas)} & \textbf{Valor total pago (R\$)} \\
\hline
1 & 8 \\
\hline
2 & 12 \\
\hline
3 & 16 \\
\hline
4 & 20 \\
\hline
\end{tabular}
\end{center}

a) Observando o padrão de variação entre o tempo de uso e o valor pago, escreva uma expressão algébrica que permita calcular o valor total V, em reais, pago por um cliente que usa a bicicleta durante t horas, para qualquer valor inteiro positivo de t.

b) Classifique o tipo de função obtida em (a), justificando sua resposta a partir da expressão algébrica encontrada.

\textbf{Gabarito proposto:} V(t) = 4t + 4; a relação é uma função polinomial do 1º grau (afim), pois apresenta taxa de variação constante (4) e t aparece com expoente 1.

\textbf{Habilidade declarada:} EM13MAT501 --- Investigar relações entre números expressos em tabelas para representá-los no plano cartesiano, identificando padrões e criando conjecturas para generalizar e expressar algebricamente essa generalização, reconhecendo quando essa representação é de função polinomial de 1º grau.
\end{samepage}

\noindent\rule{\linewidth}{0.4pt}
\textbf{Tem erro matemático?}\quad
\mbox{$\square$~não} \quad \mbox{$\square$~sim, aqui:} \underline{\hspace{5cm}}
\quad \mbox{$\square$~não sei dizer}

\textbf{Que nível de dificuldade esta questão tem para os seus alunos?}\quad
\mbox{$\square$~fácil} \quad \mbox{$\square$~média} \quad
\mbox{$\square$~difícil} \quad \mbox{$\square$~fora do alcance da turma}

\textbf{Você usaria esta questão?}\quad
\mbox{$\square$~aceita} \quad \mbox{$\square$~aceita com ajuste} \quad
\mbox{$\square$~recusada}

\textbf{Por quê?} \emph{(obrigatório quando não for ``aceita'')}

\underline{\hspace{\linewidth}}

\underline{\hspace{\linewidth}}
\vspace{0.6em}

\begin{samepage}
\subsection*{Q015}
Uma colônia de bactérias em um experimento de laboratório tem seu número de indivíduos, a partir do início da observação, dado por $N(t) = 100 \cdot 2^{t}$, em que $t$ é o tempo em horas ($t \geq 0$) e $N(t)$ é o número de bactérias. Um pesquisador quer saber, ao contrário, quanto tempo é necessário para que a colônia atinja um determinado número $N$ de bactérias, e por isso utiliza a função inversa de $N(t)$, que ele chama de $t(N)$.

a) Determine o domínio e a imagem de $N(t)$, considerando o contexto do experimento.

b) Encontre a expressão de $t(N)$ (a função inversa de $N(t)$) e determine seu domínio e sua imagem.

c) Compare o crescimento das duas funções: as duas são crescentes, mas uma cresce 'cada vez mais rápido' e a outra cresce 'cada vez mais devagar' à medida que a variável aumenta. Identifique qual é qual e justifique observando como varia o incremento de cada função.

d) Explique por que o domínio de $N(t)$ coincide com a imagem de $t(N)$, e a imagem de $N(t)$ coincide com o domínio de $t(N)$, relacionando esse fato com a posição dos gráficos das duas funções no plano cartesiano.

\textbf{Gabarito proposto:} a) $D_N=[0,+\infty)$, $Im_N=[100,+\infty)$. b) $t(N)=\log_2(N/100)$, com $D_t=[100,+\infty)$, $Im_t=[0,+\infty)$. c) $N(t)$ cresce cada vez mais rápido (convexa); $t(N)$ cresce cada vez mais devagar (côncava). d) Por serem funções inversas, seus gráficos são simétricos em relação à reta $y=x$, o que troca domínio e imagem entre as duas funções.

\textbf{Habilidade declarada:} EM13MAT403 --- Comparar e analisar as representações, em plano cartesiano, das funções exponencial e logarítmica para identificar as características fundamentais (domínio, imagem, crescimento) de cada uma, com ou sem apoio de tecnologias digitais, estabelecendo relações entre elas.
\end{samepage}

\noindent\rule{\linewidth}{0.4pt}
\textbf{Tem erro matemático?}\quad
\mbox{$\square$~não} \quad \mbox{$\square$~sim, aqui:} \underline{\hspace{5cm}}
\quad \mbox{$\square$~não sei dizer}

\textbf{Que nível de dificuldade esta questão tem para os seus alunos?}\quad
\mbox{$\square$~fácil} \quad \mbox{$\square$~média} \quad
\mbox{$\square$~difícil} \quad \mbox{$\square$~fora do alcance da turma}

\textbf{Você usaria esta questão?}\quad
\mbox{$\square$~aceita} \quad \mbox{$\square$~aceita com ajuste} \quad
\mbox{$\square$~recusada}

\textbf{Por quê?} \emph{(obrigatório quando não for ``aceita'')}

\underline{\hspace{\linewidth}}

\underline{\hspace{\linewidth}}
\vspace{0.6em}

\begin{samepage}
\subsection*{Q056}
Um biólogo monitora o crescimento de uma colônia de bactérias em um experimento de laboratório, registrando a quantidade de indivíduos ao final de cada hora, a partir da 1ª hora de observação. Os dados coletados nas quatro primeiras horas foram:

\begin{center}
\begin{tabular}{|p{0.430\linewidth}|p{0.430\linewidth}|}
\hline
\textbf{Hora (n)} & \textbf{Quantidade de bactérias} \\
\hline
1 & 100 \\
\hline
2 & 300 \\
\hline
3 & 900 \\
\hline
4 & 2700 \\
\hline
\end{tabular}
\end{center}

Essa quantidade cresce segundo o mesmo padrão observado nessas quatro medições, formando uma progressão geométrica. O biólogo deseja expressar o número de bactérias por meio de uma função exponencial $f(n)$, definida para $n$ natural, $n \geq 1$, que reproduza exatamente os valores registrados na tabela.

Assinale a alternativa que apresenta corretamente essa função $f(n)$ e o número de bactérias previsto para a 7ª hora de observação.
\begin{enumerate}[label=(\alph*),nosep,leftmargin=*]
\item $f(n) = 100 \cdot 3^{n-1}$; na 7ª hora haverá $72900$ bactérias.
\item $f(n) = 100 + 200(n-1)$; na 7ª hora haverá $1300$ bactérias.
\item $f(n) = 100 \cdot 3^{n}$; na 7ª hora haverá $218700$ bactérias.
\item $f(n) = 100 \cdot 3^{n-1}$; na 7ª hora haverá $109300$ bactérias, calculadas pela soma de todos os termos da progressão até a 7ª hora.
\end{enumerate}

\textbf{Gabarito proposto:} f(n) = 100·3\textasciicircum{}(n-1); f(7) = 72900 bactérias

\textbf{Habilidade declarada:} EM13MAT508 --- Identificar e associar sequências numéricas (PG) a funções exponenciais de domínios discretos para análise de propriedades, incluindo dedução de algumas fórmulas e resolução de problemas.
\end{samepage}

\noindent\rule{\linewidth}{0.4pt}
\textbf{Tem erro matemático?}\quad
\mbox{$\square$~não} \quad \mbox{$\square$~sim, aqui:} \underline{\hspace{5cm}}
\quad \mbox{$\square$~não sei dizer}

\textbf{Que nível de dificuldade esta questão tem para os seus alunos?}\quad
\mbox{$\square$~fácil} \quad \mbox{$\square$~média} \quad
\mbox{$\square$~difícil} \quad \mbox{$\square$~fora do alcance da turma}

\textbf{Você usaria esta questão?}\quad
\mbox{$\square$~aceita} \quad \mbox{$\square$~aceita com ajuste} \quad
\mbox{$\square$~recusada}

\textbf{Por quê?} \emph{(obrigatório quando não for ``aceita'')}

\underline{\hspace{\linewidth}}

\underline{\hspace{\linewidth}}
\vspace{0.6em}

\begin{samepage}
\subsection*{Q075}
O potencial hidrogeniônico (pH) de uma solução aquosa é definido por $pH = -\log_{10}[H^+]$, em que $[H^+]$ é a concentração de íons hidrogênio (em mol/L). Quando uma solução sofre diluição ou concentração, sua concentração de $H^+$ é multiplicada por um fator $k>0$ (por exemplo, $k = \tfrac{1}{10}$ significa que a concentração caiu para um décimo do valor original; $k=5$ significa que ela quintuplicou).

\textbf{(a)} Mostre, usando as propriedades dos logaritmos, que a variação do pH causada por esse fator, $\Delta pH = pH_{final} - pH_{inicial}$, pode ser escrita como uma fórmula que depende apenas de $k$ (isto é, não depende do valor inicial de $[H^+]$).

\textbf{(b)} Uma solução tem $pH_{inicial} = 4{,}0$. Ela é diluída de modo que a concentração de $H^+$ passa a ser $\tfrac{1}{10}$ da concentração original ($k = \tfrac{1}{10}$). Calcule o novo pH e a variação $\Delta pH$, e diga se a solução ficou mais ácida ou mais básica.

\textbf{(c)} Agora é sua vez de elaborar um problema: invente uma situação com um pH inicial diferente de $4{,}0$ e um fator $k$ diferente de $\tfrac{1}{10}$ (pode representar uma diluição, com $0<k<1$, ou uma concentração, com $k>1$). Escreva o enunciado completo do seu problema (informando claramente o pH inicial e o fator $k$ escolhidos) e, em seguida, resolva-o usando a fórmula geral obtida no item (a), calculando o novo pH, a variação $\Delta pH$ e interpretando se a solução se tornou mais ácida ou mais básica.

\textbf{Gabarito proposto:} (a) $\Delta pH = -\log_{10}(k)$, independente do pH inicial. (b) $pH_{final}=5{,}0$; $\Delta pH = +1$; a solução ficou menos ácida. (c) Resposta aberta: será considerada correta qualquer situação-problema coerente (com pH inicial e fator $k$ escolhidos pelo estudante, diferentes dos do item b) cuja resolução aplique corretamente $\Delta pH=-\log_{10}(k)$ e $pH_{final}=pH_{inicial}-\log_{10}(k)$, com interpretação correta do sinal da variação (ácida/básica).

\textbf{Habilidade declarada:} EM13MAT305 --- Resolver e elaborar problemas com funções logarítmicas nos quais é necessário compreender e interpretar a variação das grandezas envolvidas, em contextos como os de abalos sísmicos, pH, radioatividade, Matemática Financeira, entre outros.
\end{samepage}

\noindent\rule{\linewidth}{0.4pt}
\textbf{Tem erro matemático?}\quad
\mbox{$\square$~não} \quad \mbox{$\square$~sim, aqui:} \underline{\hspace{5cm}}
\quad \mbox{$\square$~não sei dizer}

\textbf{Que nível de dificuldade esta questão tem para os seus alunos?}\quad
\mbox{$\square$~fácil} \quad \mbox{$\square$~média} \quad
\mbox{$\square$~difícil} \quad \mbox{$\square$~fora do alcance da turma}

\textbf{Você usaria esta questão?}\quad
\mbox{$\square$~aceita} \quad \mbox{$\square$~aceita com ajuste} \quad
\mbox{$\square$~recusada}

\textbf{Por quê?} \emph{(obrigatório quando não for ``aceita'')}

\underline{\hspace{\linewidth}}

\underline{\hspace{\linewidth}}
\vspace{0.6em}

\begin{samepage}
\subsection*{Q034}
Duas empresas de aplicativo de transporte, A e B, calculam o valor da corrida em função da distância percorrida $x$ (em km). A empresa A cobra R\$\,2,50 por quilômetro rodado, sem nenhuma taxa adicional. A empresa B cobra uma taxa fixa de embarque de R\$\,3,00 mais R\$\,2,00 por quilômetro rodado. Sejam $y_A(x)$ e $y_B(x)$, em reais, os custos totais das corridas das empresas A e B, respectivamente, em função de $x$. Se essas duas funções forem representadas por retas em um plano cartesiano (custo $y$ no eixo vertical e distância $x$ no eixo horizontal), assinale a alternativa que descreve corretamente essas retas.
\begin{enumerate}[label=(\alph*),nosep,leftmargin=*]
\item A reta que representa a empresa A passa pela origem (0,0) e tem coeficiente angular 2,5, sendo uma função proporcional; a reta da empresa B corta o eixo y no ponto (0,3) e tem coeficiente angular 2, não sendo proporcional.
\item A reta que representa a empresa B passa pela origem (0,0), pois seu coeficiente angular é 2; a reta da empresa A corta o eixo y no ponto (0,3), pois seu coeficiente angular é 2,5.
\item Ambas as retas passam pela origem (0,0), já que representam custo por quilômetro rodado; a reta de A tem coeficiente angular 2,5 e a de B tem coeficiente angular 2.
\item A reta da empresa A corta o eixo y no ponto (0, 2,5); a reta da empresa B corta o eixo y no ponto (0,2) e tem coeficiente angular 3.
\end{enumerate}

\textbf{Gabarito proposto:} A reta que representa a empresa A passa pela origem $(0,0)$ e tem coeficiente angular 2,5 (função proporcional); a reta da empresa B corta o eixo y no ponto $(0,3)$ e tem coeficiente angular 2 (função afim, não proporcional).

\textbf{Habilidade declarada:} EM13MAT401 --- Converter representações algébricas de funções polinomiais de 1º grau para representações geométricas no plano cartesiano, distinguindo os casos nos quais o comportamento é proporcional, recorrendo ou não a softwares ou aplicativos de álgebra e geometria dinâmica.
\end{samepage}

\noindent\rule{\linewidth}{0.4pt}
\textbf{Tem erro matemático?}\quad
\mbox{$\square$~não} \quad \mbox{$\square$~sim, aqui:} \underline{\hspace{5cm}}
\quad \mbox{$\square$~não sei dizer}

\textbf{Que nível de dificuldade esta questão tem para os seus alunos?}\quad
\mbox{$\square$~fácil} \quad \mbox{$\square$~média} \quad
\mbox{$\square$~difícil} \quad \mbox{$\square$~fora do alcance da turma}

\textbf{Você usaria esta questão?}\quad
\mbox{$\square$~aceita} \quad \mbox{$\square$~aceita com ajuste} \quad
\mbox{$\square$~recusada}

\textbf{Por quê?} \emph{(obrigatório quando não for ``aceita'')}

\underline{\hspace{\linewidth}}

\underline{\hspace{\linewidth}}
\vspace{0.6em}

\begin{samepage}
\subsection*{Q002}
Uma companhia de saneamento cobra a conta de água residencial mensal com base no consumo $x$, medido em metros cúbicos ($m^3$), segundo a lei abaixo, em que $C(x)$ é o valor da conta, em reais. Considere que $x$ pode assumir qualquer valor real não negativo (não necessariamente inteiro).

$$C(x) = \begin{cases} 15 + 2x, & \text{se } 0 \le x \le 10 \\ 35 + 4(x-10), & \text{se } x > 10 \end{cases}$$

a) Determine o domínio de validade dessa função, considerando o significado de $x$ no contexto.

b) Calcule o valor da conta para um consumo de exatamente $10\,m^3$ e para um consumo de $12\,m^3$, indicando em cada caso qual das duas sentenças da lei foi utilizada.

c) Determine a imagem da função $C$.

d) A função $C$ é crescente, decrescente ou nenhuma das duas coisas em todo o seu domínio? Justifique observando o comportamento de cada sentença.

e) Descreva (sem desenhar) o formato do gráfico de $C$: de que tipo são os dois trechos, se há salto (descontinuidade) no ponto de junção $x=10$, e como se compara a inclinação de um trecho com a do outro.

\textbf{Gabarito proposto:} a) Domínio: $[0,+\infty)$. b) $C(10)=35$ (1ª sentença); $C(12)=43$ (2ª sentença). c) Imagem: $[15,+\infty)$. d) $C$ é crescente em todo o domínio. e) Duas semirretas que se encontram sem salto em $(10,35)$, sendo a segunda mais inclinada (coeficiente angular 4) que a primeira (coeficiente angular 2).

\textbf{Habilidade declarada:} EM13MAT405 --- Reconhecer funções definidas por uma ou mais sentenças (como a tabela do Imposto de Renda, contas de luz, água, gás etc.), em suas representações algébrica e gráfica, convertendo essas representações de uma para outra e identificando domínios de validade, imagem, crescimento e decrescimento.
\end{samepage}

\noindent\rule{\linewidth}{0.4pt}
\textbf{Tem erro matemático?}\quad
\mbox{$\square$~não} \quad \mbox{$\square$~sim, aqui:} \underline{\hspace{5cm}}
\quad \mbox{$\square$~não sei dizer}

\textbf{Que nível de dificuldade esta questão tem para os seus alunos?}\quad
\mbox{$\square$~fácil} \quad \mbox{$\square$~média} \quad
\mbox{$\square$~difícil} \quad \mbox{$\square$~fora do alcance da turma}

\textbf{Você usaria esta questão?}\quad
\mbox{$\square$~aceita} \quad \mbox{$\square$~aceita com ajuste} \quad
\mbox{$\square$~recusada}

\textbf{Por quê?} \emph{(obrigatório quando não for ``aceita'')}

\underline{\hspace{\linewidth}}

\underline{\hspace{\linewidth}}
\vspace{0.6em}

\begin{samepage}
\subsection*{Q078}
Em um porto, a altura da maré (em metros) varia de forma periódica ao longo do dia. Registros mostram que a maré atinge seu nível máximo de 5 m às 3h da manhã e seu nível mínimo de 1 m às 9h da manhã, repetindo esse padrão a cada 12 horas — sendo $t$ o tempo em horas contado a partir da meia-noite ($t=0$).

Um estagiário do porto, ao tentar modelar esse fenômeno, propôs a função
$$h(t) = 2\cos\left(\frac{\pi}{6}t\right) + 3,$$
com $t$ em horas, afirmando que ela representa corretamente a maré descrita.

a) Determine a amplitude, o período e o deslocamento vertical do fenômeno periódico descrito pelos registros.

b) Escreva a função $h(t)$, na forma $h(t) = A\cos(B(t-C)) + D$, que representa corretamente a maré observada.

c) Avalie se a proposta do estagiário está correta, comparando os valores previstos por sua função com os dados observados nos instantes $t=3$ h e $t=9$ h. Caso ela esteja incorreta, explique qual característica do gráfico (amplitude, período ou deslocamento de fase) foi tratada de forma equivocada.

\textbf{Gabarito proposto:} a) Amplitude $A=2$ m, período $T=12$ h, deslocamento vertical $D=3$ m. b) $h(t) = 2\cos\left(\frac{\pi}{6}(t-3)\right)+3$. c) A proposta do estagiário está incorreta: amplitude e período estão certos, mas falta o deslocamento de fase ($C=3$); sua função dá $h(3)=3$ e $h(9)=3$, valores que não correspondem aos dados reais (5 m e 1 m).

\textbf{Habilidade declarada:} EM13MAT306 --- Resolver e elaborar problemas em contextos que envolvem fenômenos periódicos reais, como ondas sonoras, ciclos menstruais, movimentos cíclicos, entre outros, e comparar suas representações com as funções seno e cosseno, no plano cartesiano, com ou sem apoio de aplicativos de álgebra e geometria.
\end{samepage}

\noindent\rule{\linewidth}{0.4pt}
\textbf{Tem erro matemático?}\quad
\mbox{$\square$~não} \quad \mbox{$\square$~sim, aqui:} \underline{\hspace{5cm}}
\quad \mbox{$\square$~não sei dizer}

\textbf{Que nível de dificuldade esta questão tem para os seus alunos?}\quad
\mbox{$\square$~fácil} \quad \mbox{$\square$~média} \quad
\mbox{$\square$~difícil} \quad \mbox{$\square$~fora do alcance da turma}

\textbf{Você usaria esta questão?}\quad
\mbox{$\square$~aceita} \quad \mbox{$\square$~aceita com ajuste} \quad
\mbox{$\square$~recusada}

\textbf{Por quê?} \emph{(obrigatório quando não for ``aceita'')}

\underline{\hspace{\linewidth}}

\underline{\hspace{\linewidth}}
\vspace{0.6em}

\begin{samepage}
\subsection*{Q047}
Considere as três funções quadráticas, definidas para todo $x \in \mathbb{R}$:

$f(x) = 3x^2$

$g(x) = 3x^2 - 12$

$h(x) = 3(x+2)^2$

a) Entre as três funções, identifique qual delas representa uma situação em que a variável $y$ é diretamente proporcional ao quadrado de $x$, isto é, pode ser escrita na forma $y = kx^2$ com $k$ constante. Justifique algebricamente por que as outras duas NÃO se enquadram nessa relação de proporcionalidade direta.

b) Determine as coordenadas do vértice da parábola associada a cada uma das três funções.

c) Descreva, em termos de translações no plano cartesiano (direção e número de unidades deslocadas), como os gráficos de $g$ e de $h$ podem ser obtidos a partir do gráfico de $f$.

\textbf{Gabarito proposto:} Apenas $f(x)=3x^2$ é diretamente proporcional a $x^2$ (vértice em $(0,0)$). Vértices: $f\to(0,0)$, $g\to(0,-12)$, $h\to(-2,0)$. O gráfico de $g$ é o de $f$ deslocado 12 unidades para baixo; o de $h$ é o de $f$ deslocado 2 unidades para a esquerda.

\textbf{Habilidade declarada:} EM13MAT402 --- Converter representações algébricas de funções polinomiais de 2º grau para representações geométricas no plano cartesiano, distinguindo os casos nos quais uma variável for diretamente proporcional ao quadrado da outra, recorrendo ou não a softwares ou aplicativos de álgebra e geometria dinâmica.
\end{samepage}

\noindent\rule{\linewidth}{0.4pt}
\textbf{Tem erro matemático?}\quad
\mbox{$\square$~não} \quad \mbox{$\square$~sim, aqui:} \underline{\hspace{5cm}}
\quad \mbox{$\square$~não sei dizer}

\textbf{Que nível de dificuldade esta questão tem para os seus alunos?}\quad
\mbox{$\square$~fácil} \quad \mbox{$\square$~média} \quad
\mbox{$\square$~difícil} \quad \mbox{$\square$~fora do alcance da turma}

\textbf{Você usaria esta questão?}\quad
\mbox{$\square$~aceita} \quad \mbox{$\square$~aceita com ajuste} \quad
\mbox{$\square$~recusada}

\textbf{Por quê?} \emph{(obrigatório quando não for ``aceita'')}

\underline{\hspace{\linewidth}}

\underline{\hspace{\linewidth}}
\vspace{0.6em}

\begin{samepage}
\subsection*{Q022}
Uma empresa de entregas rápidas registrou o valor cobrado (em reais) para diferentes distâncias percorridas por seus motoboys, conforme a tabela abaixo:

\begin{center}
\begin{tabular}{|p{0.430\linewidth}|p{0.430\linewidth}|}
\hline
\textbf{Distância d (km)} & \textbf{Valor cobrado V (R\$)} \\
\hline
2 & 12 \\
\hline
5 & 21 \\
\hline
8 & 30 \\
\hline
11 & 39 \\
\hline
\end{tabular}
\end{center}

Analisando os dados da tabela, qual expressão algébrica V(d) generaliza corretamente o padrão observado, e que tipo de função ela representa?
\begin{enumerate}[label=(\alph*),nosep,leftmargin=*]
\item V(d) = 3d + 6 — função polinomial do 1º grau
\item V(d) = 6d — função polinomial do 1º grau
\item V(d) = 3d² + 6 — função polinomial do 2º grau
\item V(d) = 9d - 6 — função polinomial do 1º grau
\end{enumerate}

\textbf{Gabarito proposto:} V(d) = 3d + 6 — função polinomial do 1º grau (afim)

\textbf{Habilidade declarada:} EM13MAT501 --- Investigar relações entre números expressos em tabelas para representá-los no plano cartesiano, identificando padrões e criando conjecturas para generalizar e expressar algebricamente essa generalização, reconhecendo quando essa representação é de função polinomial de 1º grau.
\end{samepage}

\noindent\rule{\linewidth}{0.4pt}
\textbf{Tem erro matemático?}\quad
\mbox{$\square$~não} \quad \mbox{$\square$~sim, aqui:} \underline{\hspace{5cm}}
\quad \mbox{$\square$~não sei dizer}

\textbf{Que nível de dificuldade esta questão tem para os seus alunos?}\quad
\mbox{$\square$~fácil} \quad \mbox{$\square$~média} \quad
\mbox{$\square$~difícil} \quad \mbox{$\square$~fora do alcance da turma}

\textbf{Você usaria esta questão?}\quad
\mbox{$\square$~aceita} \quad \mbox{$\square$~aceita com ajuste} \quad
\mbox{$\square$~recusada}

\textbf{Por quê?} \emph{(obrigatório quando não for ``aceita'')}

\underline{\hspace{\linewidth}}

\underline{\hspace{\linewidth}}
\vspace{0.6em}

\begin{samepage}
\subsection*{Q040}
Duas empresas de transporte por aplicativo, Moto Rápida e Carro Center, cobram uma corrida com base na distância percorrida. A tabela abaixo mostra o valor cobrado, em reais, para algumas distâncias, em quilômetros:

Moto Rápida: 2 km $\to$ R\$\,7,00; 5 km $\to$ R\$\,17,50; 8 km $\to$ R\$\,28,00

Carro Center: 2 km $\to$ R\$\,9,00; 5 km $\to$ R\$\,16,50; 8 km $\to$ R\$\,24,00

Um estudante decide representar, em um mesmo plano cartesiano, o custo da corrida (eixo y, em reais) em função da distância percorrida (eixo x, em km) para cada uma das duas empresas, obtendo duas retas.

Assinale a alternativa que descreve corretamente essas duas retas, indicando qual delas representa uma relação de proporcionalidade direta entre custo e distância, e em que ponto (distância, custo) as duas retas se cruzam.
\begin{enumerate}[label=(\alph*),nosep,leftmargin=*]
\item A reta de Moto Rápida passa pela origem, com inclinação 3,5, representando proporcionalidade direta; a reta de Carro Center corta o eixo y em (0; 4) e tem inclinação 2,5; as retas se cruzam no ponto (4; 14).
\item A reta de Carro Center passa pela origem, com inclinação 2,5, representando proporcionalidade direta; a reta de Moto Rápida corta o eixo y em (0; 4) e tem inclinação 3,5; as retas se cruzam no ponto (4; 14).
\item A reta de Moto Rápida passa pela origem, com inclinação 3,5, representando proporcionalidade direta; a reta de Carro Center corta o eixo y em (0; 4) e tem inclinação 2,5; as retas se cruzam no ponto (4; 10).
\item Tanto a reta de Moto Rápida (inclinação 3,5) quanto a de Carro Center (inclinação 2,5) passam pela origem, sendo ambas proporcionais; por isso, as retas só têm um ponto em comum, a origem (0; 0).
\end{enumerate}

\textbf{Gabarito proposto:} A reta de Moto Rápida passa pela origem, com inclinação 3,5, representando proporcionalidade direta; a reta de Carro Center corta o eixo y em (0; 4) e tem inclinação 2,5; as retas se cruzam no ponto (4; 14).

\textbf{Habilidade declarada:} EM13MAT401 --- Converter representações algébricas de funções polinomiais de 1º grau para representações geométricas no plano cartesiano, distinguindo os casos nos quais o comportamento é proporcional, recorrendo ou não a softwares ou aplicativos de álgebra e geometria dinâmica.
\end{samepage}

\noindent\rule{\linewidth}{0.4pt}
\textbf{Tem erro matemático?}\quad
\mbox{$\square$~não} \quad \mbox{$\square$~sim, aqui:} \underline{\hspace{5cm}}
\quad \mbox{$\square$~não sei dizer}

\textbf{Que nível de dificuldade esta questão tem para os seus alunos?}\quad
\mbox{$\square$~fácil} \quad \mbox{$\square$~média} \quad
\mbox{$\square$~difícil} \quad \mbox{$\square$~fora do alcance da turma}

\textbf{Você usaria esta questão?}\quad
\mbox{$\square$~aceita} \quad \mbox{$\square$~aceita com ajuste} \quad
\mbox{$\square$~recusada}

\textbf{Por quê?} \emph{(obrigatório quando não for ``aceita'')}

\underline{\hspace{\linewidth}}

\underline{\hspace{\linewidth}}
\vspace{0.6em}

\begin{samepage}
\subsection*{Q049}
Cientistas utilizam escalas logarítmicas para descrever a magnitude de fenômenos naturais cuja intensidade varia em muitas ordens de grandeza, como ocorre com a energia liberada em abalos sísmicos. Um modelo simplificado dessa relação define a magnitude $M$ de um terremoto, em função da energia $E$ (em joules) por ele liberada, pela expressão

$$M(E) = \frac{2}{3}\left(\log_{10}(E) - 4\right), \quad E > 0.$$

a) Um terremoto liberou energia $E_1 = 10^{13}$ J. Calcule sua magnitude $M_1$.

b) Um segundo terremoto ocorreu na mesma região, com magnitude exatamente uma unidade maior que a do primeiro, ou seja, $M_2 = M_1 + 1$. Determine, em função de $E_1$, a energia $E_2$ liberada por esse segundo terremoto, calcule a razão $E_2/E_1$ e mostre que essa razão não depende do valor de $E_1$ — isto é, que qualquer terremoto cuja magnitude seja uma unidade maior que a de outro libera sempre o mesmo número de vezes mais energia, independentemente do nível de energia inicial.

\textbf{Gabarito proposto:} M1 = 6; E2/E1 = 10\textasciicircum{}(3/2) = 10$\surd$10 $\approx$ 31,62, valor constante e independente de E1 (todo acréscimo de 1 unidade na magnitude multiplica a energia por esse mesmo fator).

\textbf{Habilidade declarada:} EM13MAT305 --- Resolver e elaborar problemas com funções logarítmicas nos quais é necessário compreender e interpretar a variação das grandezas envolvidas, em contextos como os de abalos sísmicos, pH, radioatividade, Matemática Financeira, entre outros.
\end{samepage}

\noindent\rule{\linewidth}{0.4pt}
\textbf{Tem erro matemático?}\quad
\mbox{$\square$~não} \quad \mbox{$\square$~sim, aqui:} \underline{\hspace{5cm}}
\quad \mbox{$\square$~não sei dizer}

\textbf{Que nível de dificuldade esta questão tem para os seus alunos?}\quad
\mbox{$\square$~fácil} \quad \mbox{$\square$~média} \quad
\mbox{$\square$~difícil} \quad \mbox{$\square$~fora do alcance da turma}

\textbf{Você usaria esta questão?}\quad
\mbox{$\square$~aceita} \quad \mbox{$\square$~aceita com ajuste} \quad
\mbox{$\square$~recusada}

\textbf{Por quê?} \emph{(obrigatório quando não for ``aceita'')}

\underline{\hspace{\linewidth}}

\underline{\hspace{\linewidth}}
\vspace{0.6em}

\begin{samepage}
\subsection*{Q088}
Um biólogo monitora duas culturas de microrganismos mantidas em condições ideais de crescimento.

\begin{itemize}[nosep,leftmargin=*]
\item A cultura A começa com uma quantidade $A_0$ de microrganismos, e essa quantidade dobra a cada 2 horas.
\item A cultura B começa com o quádruplo da quantidade inicial da cultura A, ou seja, com $4A_0$ microrganismos, e essa quantidade dobra a cada 6 horas.
\end{itemize}

Sabendo que as duas populações crescem de forma exponencial desde o início da observação, determine depois de quantas horas as duas culturas terão exatamente a mesma quantidade de microrganismos.
\begin{enumerate}[label=(\alph*),nosep,leftmargin=*]
\item 6 horas
\item 3 horas
\item 8 horas
\item 4 horas
\end{enumerate}

\textbf{Gabarito proposto:} 6 horas

\textbf{Habilidade declarada:} EM13MAT304 --- Resolver e elaborar problemas com funções exponenciais nos quais é necessário compreender e interpretar a variação das grandezas envolvidas, em contextos como o da Matemática Financeira e o do crescimento de seres vivos microscópicos, entre outros.
\end{samepage}

\noindent\rule{\linewidth}{0.4pt}
\textbf{Tem erro matemático?}\quad
\mbox{$\square$~não} \quad \mbox{$\square$~sim, aqui:} \underline{\hspace{5cm}}
\quad \mbox{$\square$~não sei dizer}

\textbf{Que nível de dificuldade esta questão tem para os seus alunos?}\quad
\mbox{$\square$~fácil} \quad \mbox{$\square$~média} \quad
\mbox{$\square$~difícil} \quad \mbox{$\square$~fora do alcance da turma}

\textbf{Você usaria esta questão?}\quad
\mbox{$\square$~aceita} \quad \mbox{$\square$~aceita com ajuste} \quad
\mbox{$\square$~recusada}

\textbf{Por quê?} \emph{(obrigatório quando não for ``aceita'')}

\underline{\hspace{\linewidth}}

\underline{\hspace{\linewidth}}
\vspace{0.6em}

\newpage
\subsection*{Para terminar}

\textbf{1. Você usaria um sistema assim no seu planejamento? Por quê?}

\underline{\hspace{\linewidth}}

\underline{\hspace{\linewidth}}

\bigskip
\textbf{2. O que faltou nas questões que você viu?}

\underline{\hspace{\linewidth}}

\underline{\hspace{\linewidth}}

\bigskip
\textbf{3. O que atrapalhou --- na questão, no enunciado, no formato deste
material?}

\underline{\hspace{\linewidth}}

\underline{\hspace{\linewidth}}

\bigskip
Obrigado. Se quiser saber quais itens tinham erro e qual era, é só pedir.
\end{document}